\section{Архитектура олимпиады}

\begin{frame}{Четыре этапа}
  \vspace{0.2em}
  \small
  \begin{tabular}{@{}p{0.14\textwidth}p{0.17\textwidth}p{0.16\textwidth}p{0.21\textwidth}p{0.24\textwidth}@{}}
    \toprule
    \textbf{Этап} & \textbf{Уровень} & \textbf{Формат} & \textbf{Длительность} & \textbf{Кто обеспечивает}\\
    \midrule
    I   & Школьный      & Дистанционно & 2 часа, скользящее окно старта & Участник / учреждение образования \\
    \addlinespace[0.3em]
    II  & Районный      & Дистанционно & 2 часа, скользящее окно старта & Участник / отдел (управление) образования \\
    \addlinespace[0.3em]
    III & Областной     & \textbf{Очно} & 2 тура по 4--6 часов, в разные дни & \textbf{Оргкомитеты III этапа} \\
    \addlinespace[0.3em]
    IV  & Заключительный& \textbf{Очно}, Технопарк & 2 тура по 4--6 часов, в разные дни & Республиканский Оргкомитет \\
    \bottomrule
  \end{tabular}

  \vspace{0.9em}
  {\footnotesize\color{acGray}
  На I и II этапах применяется прокторинг: запись экрана, камеры, микрофона и аудиовыхода.
  Единый набор заданий на III и IV этапах обязателен, время старта туров фиксировано.}
\end{frame}
